% FDUletter_example.tex - an example latex file to illustrate FDUletter.cls
%
% Template by Brian Wood (brian.wood@oregonstate.edu).  Please feel free to send suggestions for changes; this template/cls is not exactly elegantly done!
% Modified by Huang Weiran (huangweiran1998@outlook.com) to fit the need of FDU students.


\documentclass[12pt]{OSUletter}
\usepackage{fontspec} 
\usepackage{tikz} 
\usepackage{xcolor}
	\definecolor{osuorange}{rgb}{0.0549, 0.208, 0.545}
	\definecolor{osuorange}{HTML}{D73F09}
	\definecolor{osuwhite}{HTML}{FFFFFF}
	\definecolor{osublack}{HTML}{000000}
\usepackage{lipsum}
\usepackage{fancyhdr}
\usepackage{lastpage}
\usepackage{eso-pic}
%
% This section is just a bunch of busywork so that the second and following pages read ``Page X of Y''
\pagestyle{fancy}
\fancyhf{}
\renewcommand{\headrulewidth}{0pt}
\renewcommand{\footrulewidth}{0pt}
\rhead{Page \thepage \hspace{1pt} of \pageref{LastPage}}
%
%
% Set custom font here. Comment this line out if you do not have a Cambria font (originally included with this template) installed; computer modern (or whatever your current default font is) will be substituted.
%
%\setmainfont{[Cambria.ttf]}[BoldFont  = [CambriaBold.ttf], ItalicFont  = [CambriaItalic.ttf], BoldItalicFont = [CambriaBoldItalic.ttf] ]

\newcommand{\watermark}[3]{\AddToShipoutPictureBG{
\parbox[b][\paperheight]{\paperwidth}{
\vfill%
\centering%
\begin{tikzpicture}
    \path (0,0) -- (\paperwidth,\paperheight);
    \node[opacity=.06] at (current page.center)
    {\includegraphics[width=0.75\textwidth]{cover-letter/osu-background.png}};
    \end{tikzpicture}
\vfill}}}

% The material below is a whole big dang thing whose purpose is just to set up a fixed coordinate system for \tikz so that you can put the Department or School address in the upper right-hand side without it moving all around every time you change something in the page.  I think it works.
% Defining a new coordinate system for the page:
%
% --------------------------
% |(-1,1)    (0,1)    (1,1)|
% |                        |
% |(-1,0)    (0,0)    (1,0)|
% |                        |
% |(-1,-1)   (0,-1)  (1,-1)|
% --------------------------
\makeatletter
\def\parsecomma#1,#2\endparsecomma{\def\page@x{#1}\def\page@y{#2}}
\tikzdeclarecoordinatesystem{page}{
    \parsecomma#1\endparsecomma
    \pgfpointanchor{current page}{north east}
    % Save the upper right corner
    \pgf@xc=\pgf@x%
    \pgf@yc=\pgf@y%
    % save the lower left corner
    \pgfpointanchor{current page}{south west}
    \pgf@xb=\pgf@x%
    \pgf@yb=\pgf@y%
    % Transform to the correct placement
    \pgfmathparse{(\pgf@xc-\pgf@xb)/2.*\page@x+(\pgf@xc+\pgf@xb)/2.}
    \expandafter\pgf@x\expandafter=\pgfmathresult pt
    \pgfmathparse{(\pgf@yc-\pgf@yb)/2.*\page@y+(\pgf@yc+\pgf@yb)/2.}
    \expandafter\pgf@y\expandafter=\pgfmathresult pt
}
\makeatother
%
%
%%%%%%%%%%% Put Personal Information Here %%%%%%%%%%%
%
\def\name{Anthony Walker}
%
% List your degree(s), licences, etc. here if you like.
%\def\What{, Your degrees, etc.} 
%
% Set the name of your Department or School here
% I honestly don't know why the negative spacing is necessary to get the alignment of the first line correct.  This must be a ``\tikz'' thing.
%%%%%%%%%%%%%%%%%%  School or Department %%%%%%%%%%%%%%%
\def\Where{\hspace{-1.2mm}\textbf{\color{osuorange}%{osuorange}
School of MIME,\\ Oregon State University
}} 
%%%%%%%%%%%%  Additional Contact Information %%%%%%%%%%%
%
% Set your preferred primary contact address here.
%
\def\CityZip{Corvallis, 97330} 
%
% Set your FDU e-mail here
\def\Email{\textbf{\color{osuorange}E-mail}: walkanth@oregonstate.edu}
%
% Set your preferred contact number here
\def\TEL{\textbf{\color{osuorange}Phone}: 707-337-3595}
%
% % Set your department or personal website here
% \def\URL{\textbf{\color{osuorange}URL}: http://www.fudan.edu.cn/}
%
%%%%%%%%%%%%%%%%  Footer information  %%%%%%%%%%%%%%%%%%
%
%  The next line is for your college, used as a footer.  If you prefer not to have this, just comment out these lines in favor of the line labeled "[[Alternate]]" below
% \def\school{\small{
%   FDU $\cdot$
%      ~School of Computer Science $\cdot$
%      ~No.220 Handan RD $\cdot$
%      ~Shanghai, China P.R} } 
% \def\school{~}  % [[Alternate]]
%
%%%%%%%%%%%%%%%%%%%%%  Signature line  %%%%%%%%%%%%%%%%%%%%%
%
% Set your signature line here.
% One can add a signature image in a PDF file using the following code; this requires a file called "signature_block.pdf" to be installed in the same folder as the .tex file.  The vertical spacing (\vspace) and the scaling will have to be adjusted to get things to look correct for your particular signature image. Alternatively, comment out the following line in favor of the one labeled "[[Alternate]]" if you want to sign a paper copy of the letter.
%
\signature{ 
\vspace{-12mm}\includegraphics[]{cover-letter/signature.png}\\\vspace{-2mm}
\name}
%\signature{\name}  % [[Alternate]]

% This block sets up the address on the right-hand side of the header. 
%
% The following lines just compile the information you set up into the LaTex letter variable "address" for later use.
%
%The following command "clears out" the default address so that it can be better set using \tikz
\address{}

\def\newaddress{
\Where\\ 
\Address\\ 
\CityZip\\ 
\TEL\\ 
\Email\\ 
\URL 
}

%%%%%%%%%%%  DATE  %%%%%%%%%%%%%%%%%%%%%%%%%
% If you want a date different from the current date, comment out the next line in favor of the line labeled "[[Alternate]]".  The ``\vspace{10mm}'' just moves the date down a tiny bit so it doesn't interfere with the header.  This can be adjusted to your preference.
%
\date{\vspace{10mm} \today}
%\date{\vspace{10mm} 20 September 2020}  %[[Alternate]]
%
%%%%%%%%%%% Set the subject here if there is one  %%%%
%\subject{Stuff} % optional subject line

\newcommand{\position}{Operations Engineer, Intern}

\begin{document}
%
%
%%%%%%%%  The "To" address goes here.
%
% \begin{letter}{
%               Some University\\ 
%               Some Addresss\\ 
%               SomeTown, SomeState 					       				  		 ~~SomeZip
%               }
% This line sets up the return address to the right-side of the OSU logo.  The location is set with absolute node addresses using ``\tikz''.  It can still be a bit fussy, and you may need to alter this a little to get things to look right.  The bit that changes the position are the numbers in parentheses ``at (14.2,2.7)''
%
\begin{tikzpicture}[remember picture,overlay,,every node/.style={anchor=center}]
\node[text width=7cm] at (page cs:0.5,0.73){\small \newaddress};
\end{tikzpicture} 

%%%%%%  The ``opening'' is just the method of address you would like to use at the start of the letter.
%
\opening{To Whom It May Concern,}

%%%%%%%%%% Body of letter   %%%%%%%%%%%%%%
% Remove it if you do not want watermark
\watermark{}{}{}

My name is Anthony Walker and I am a fourth year PhD student at Oregon State University and I am interested in the \position position. I am qualified for this position because I have a wide breadth of experience in computer science and software development.
Nominally, my degree is in mechanical engineering but my work is a combination of computer science, software development, physics, chemistry, mathematics, etc. I also have substantial course work and minors in computer science.

I have software related experience in both teaching and research. I have served as a teaching assistant for multiple classes requiring programming and I am currently scheduled to teach a Python course at Oregon State University in the winter 2021 term. My research also provides me with software experience as it requires that I implement new functionality in an open-source chemical kinetics package called Cantera. My day to day revolves mostly around the use of C++, Python, and Bash but I have used many languages which can be found on my resume, ranked by experience. Similarly, my Master's project was the implementation of heterogeneous communication avoidance algorithms for accelerating simulations on distributed systems. To this end, I used a number of languages and software packages listed on my resume. Finally, I have sought personal projects/tutorials which have extended my skill set into various fields such as cybersecurity, application development, and web development. 

My skills aside, I believe my outlook aligns well with the Duolingo's mission. I am an open-source developer for Cantera and believe that software and educational material should be available for everyone without a subtantial expense. My experience with teaching has also given me an appreciation for the learning process of individuals. I have faced multiple challenges developing "the correct way" to teach a specific student which is always a learning experience for me.






%%%%%%% ``closing'' sets the sign-off line.
\closing{Sincerely,}

% Comment out/in the lines below as necessary
%\encl{If an enclosure is provided, let them know what it is.}

%\ps{A postscript if that is a thing you do.}

%\cc{Someone Who Cares (and is copied).}

\end{letter}

\end{document}